\chapter{Weighted Charges and Ground-State Fidelity}
\label{ch:weighted-fidelity}
The purpose of the auxiliary operator is to turn a bound on the
weighted charge energy of a trial state into information about its
projection onto the original kernel. Work at fixed $\varepsilon>0$,
with the positive radial weight
$w_\varepsilon(X)=(\|X\|^2+\varepsilon^2)^{1/2}$ and the bounded
reciprocal output multiplier $B_\varepsilon=w_\varepsilon^{-1}$.
The weighted charge graph consists of pairs $(u,z)$ for which
$(u,B_\varepsilon z)$ belongs to the original charge graph. This route
uses continuity and injectivity of the reciprocal multiplier to obtain
a closed, single-valued graph. The source domain and density proofs
establish the hypotheses for forming its adjoint square.

\section{The weighted Hamiltonian and the shared kernel}
Let $S_\varepsilon$ be the accepted weighted charge column on the
original physical input space. Its normalized operator and form are
\[
 H_\varepsilon=\tfrac1{16}S_\varepsilon^\dagger S_\varepsilon,
 \qquad q_\varepsilon[u]=\tfrac1{16}\|S_\varepsilon u\|^2 .
\]
The accepted weak form uses the natural domain of this adjoint square.
For either charge column, the kernel of its adjoint square equals
the kernel of the column. On the weighted graph, $(u,0)$ belongs
precisely when $(u,B_\varepsilon 0)=(u,0)$ belongs to the original
graph. This is the mechanism behind
\begin{equation}\label{eq:ed08-kernels}
 \ker H_\varepsilon=\ker H\qquad(\varepsilon>0).
\end{equation}
The equality concerns subspaces of the original physical Hilbert
space. The common kernel may be trivial. The zero-regulator limit is
a different analysis.
\formalrecord{The natural adjoint square, normalized weak form, and
kernel comparison have separate receipts in atlas entries
\texttt{E08-C4A}, \texttt{E08-C4B}, and \texttt{E08-C4C}. The
reciprocal pairing is indexed by \texttt{E08-CERTP0}.}

\section{A conditional fidelity bound}
\label{sec:ed08-resolvent}
Let $E$ be the original physical Hilbert space and $T$ its closed
sixteen-charge column, with $H=T^\dagger T/16$. For a normalized
physical singlet $\psi$ in the weighted domain put
$z=S_\varepsilon\psi$ and $\chi_\alpha=(T\psi)_\alpha$.
The existing reciprocal construction gives
$\chi_\alpha=B_\varepsilon z_\alpha$ under its actual graph
hypotheses. Assume additionally that the outputs are physical,
belong to a reducing spinor sector, and satisfy, for $t>0$,
\begin{equation}\label{eq:ed08-R}
 \tfrac1{16}\sum_\alpha
 \langle\chi_\alpha,(H+t)^{-1}\chi_\alpha\rangle
 =\langle\psi,H(H+t)^{-1}\psi\rangle.
\end{equation}
Assume the native spinor form-domain inequality
\begin{equation}\label{eq:ed08-Hardy}
 \|B_\varepsilon\eta\|^2\leq\kappa h[\eta],
 \qquad \eta\in D(h)\cap E_{\rm spinor}.
\end{equation}
Neither displayed assumption has been qualified on the original
source. The constant is not imported from a paper by convention.

Conditionally, take $\eta_\alpha=(H+t)^{-1}\chi_\alpha$ and
$a_\alpha=\langle\chi_\alpha,\eta_\alpha\rangle\geq0$.
The resolvent equation gives $h[\eta_\alpha]\leq a_\alpha$.
Reciprocal pairing and Cauchy--Schwarz yield
$a_\alpha\leq\sqrt\kappa\|z_\alpha\|\sqrt{a_\alpha}$,
hence $a_\alpha\leq\kappa\|z_\alpha\|^2$. Averaging and
letting the spectral shift $t$ decrease to zero gives
$\|(I-P_{\ker H})\psi\|^2\leq\kappa q_\varepsilon[\psi]$.
The radial regulator $\varepsilon$ stays fixed. This is a conditional
argument using bounded shifted resolvents, not a new Lean certificate
or a proof of the assumptions.


\section{What would turn the bound into a certificate}
For a normalized original physical trial vector, the conditional
conclusion is
\[
 \|P_{\ker H}\psi\|^2\geq\max\{0,1-\kappa q_\varepsilon[\psi]\}.
\]
A strictly positive right side would establish a nonzero vector in
the original ground subspace. The work still requires a qualified
physical singlet trial vector in the weighted form domain, a
source-realized reducing spinor sector containing its charge outputs,
the averaged resolvent identity, and a Hardy estimate with a usable
constant. Only then can the trial energy be compared to $1/\kappa$.
The identity $H=T^\dagger T/16$ supplies the closed positive form.
A square identity for each individual charge or an intertwiner through
the resolvent requires its own domain and covariance argument.
